% Part: sets-functions-relations
% Chapter: arithmetization
% Section: checking-details
%
\documentclass[../../../include/open-logic-section]{subfiles}


\begin{document}
	
\olfileid[fa-IR]{sfr}{arith}{check}
\olsection{حلقه‌ها و میدان‌های مرتب}

در سراسر این فصل ادعا کردیم که تعریف‌هایی معین «چنان‌که باید» رفتار
می‌کنند. در این پیوست فنی، مقصود خود را به‌تفصیل بیان می‌کنیم و (طرحی از
این را) نشان می‌دهیم که تعریف‌ها چگونه «به‌درستی» رفتار می‌کنند.

در \olref[int]{sec} جمع و ضرب را روی $\Int$ تعریف کردیم. می‌خواهیم نشان
دهیم که این تعریف‌ها ساختاری را که از آن‌ها «می‌خواهیم» به $\Int$ می‌دهند.
به‌ویژه، ساختار مورد نظر یک حلقهٔ جابجایی است.

\begin{defn}
	یک \emph{حلقهٔ جابجایی} مجموعه‌ای چون $S$ است که به اعضای مشخص $0$ و
	$1$ و اعمال $+$ و $\times$ مجهز است و این هشت فرمول را برآورده می‌کند:
	\begin{align*}
		\emph{شرکت‌پذیری}&&a + (b+ c) & = (a + b) + c \\
		&& (a \times b) \times c & = a \times (b\times c)\\
		\emph{جابجایی}&&a + b &= b+ a  \\
		&&  a \times b&= b\times a\\ 
		\emph{همانی‌ها}&&a + 0 &= a \\
		&& a \times 1 &= a\\
		\emph{وارون جمعی}&&(\exists b\in S)0&=a + b\\
		\emph{توزیع‌پذیری}&&a \times (b+ c ) &= (a \times b) + (a \times c)
	\end{align*}
	به‌طور ضمنی، همهٔ این فرمول‌ها با کمیت‌گذارهای کلیِ محدود به $S$ بسته
	شده‌اند. همچنین توجه کنید که اعضای $0$ و~$1$ در اینجا لازم نیست همان
	اعداد طبیعیِ هم‌نام باشند.
\end{defn}

پس برای بررسی اینکه اعداد صحیح حلقه‌ای جابجایی تشکیل می‌دهند، تنها باید
بررسی کنیم که این هشت شرط را برآورده می‌کنیم. اثبات هیچ‌یک از شرط‌ها
{دشوار} نیست، اما اندکی پرزحمت است. برای نمونه، در حالت جمع،
\emph{شرکت‌پذیری} را چنین ثابت می‌کنیم:

\begin{proof} 
$i, j, k \in \Int$ را ثابت بگیرید. پس $a_1, b_1, a_2, b_2, a_3, b_3 \in
\Nat$ وجود دارند چنان‌که $i = \equivrep{a_1, b_1}{}$ و $j =
\equivrep{a_2,b_2}{}$ و $k = \equivrep{a_3, b_3}{}$. (برای خوانایی،
«$\equivrep{x, y}{}$» را به‌جای «$\equivrep{x, y}{\Intequiv}$» می‌نویسیم؛ در سراسر
این بخش نیز چنین خواهیم کرد.) اکنون:
\begin{align*}
	i + (j + k) &= \equivrep{a_1, b_1}{}+(\equivrep{a_2, b_2}{} + \equivrep{a_3, b_3}{}) \\
	&= \equivrep{a_1,  b_1}{} + \equivrep{a_2+a_3, b_2+b_3}{}\\
	&= \equivrep{a_1 + (a_2 + a_3), b_1 + (b_2 + b_3)}{}\\
	&= \equivrep{(a_1 + a_2) + a_3, (b_1 + b_2) + b_3}{}\\
	&= \equivrep{a_1 + a_2, b_1 + b_2}{} + \equivrep{a_3, b_3}{}\\
	&= (\equivrep{a_1, b_1}{} + \equivrep{a_2, b_2}{}) + \equivrep{a_3, b_3}{}\\
	&= (i+j) + k
\end{align*}
و در این کار آزادانه از رفتار جمع روی $\Nat$ بهره می‌گیریم.
\end{proof}

به همین ترتیب، \emph{وارون جمعی} را چنین ثابت می‌کنیم:

\begin{proof}
$i \in \Int$ را ثابت بگیرید؛ پس $i = \equivrep{a,b}{}$ برای بعضی $a,b \in
\Nat$. بگذارید $j = \equivrep{b,a}{} \in \Int$. با بهره‌گیری
از رفتار اعداد طبیعی، $(a+b) + 0 = 0 + (a+b)$، و بنابراین بنا به تعریف
$\tuple{a+b, b+a} \sim_\Int \tuple{0,0}$، و ازاین‌رو
$\equivrep{a+b, b+a}{} = \equivrep{0, 0}{} = 0_\Int$. اکنون $i + j =
\equivrep{a,b}{}+\equivrep{b,a}{}=\equivrep{a+b, b+a}{}= \equivrep{0,
0}{} = 0_\Int$.
\end{proof}

و این هم برهانی برای \emph{توزیع‌پذیری}:

\begin{proof}
مانند بالا، $i = \equivrep{a_1, b_1}{}$ و $j =
\equivrep{a_2,b_2}{}$ و $k = \equivrep{a_3, b_3}{}$ را ثابت بگیرید. اکنون:
\begin{align*}
	i \times (j + k) 
	&= \equivrep{a_1, b_1}{} \times (\equivrep{a_2,b_2}{} + \equivrep{a_3, b_3}{})\\
	&= \equivrep{a_1, b_1}{} \times \equivrep{a_2 + a_3,b_2+b_3}{}\\
	&= \equivrep{a_1  (a_2 + a_3) + b_1  (b_2+b_3), a_1  (b_2 + b_3) + b_1 (a_2 + a_3)}{}\\
	&= \equivrep{a_1 a_2 + a_1a_3 + b_1 b_2+b_1b_3, a_1 b_2 + a_1b_3 + a_2b_1 + a_3b_1}{}\\		
%		&= \equivrep{a_1 a_2 + b_1b_2 + a_1a_3 + b_1 b_2+b_1b_3, a_1 b_2 + a_1b_3 + a_2b_1 + a_3b_1}{}\\		
	&= \equivrep{a_1a_2 + b_1b_2, a_1b_2 + a_2b_1}{} + \equivrep{a_1a_3 + b_1b_3, a_1b_3 + a_3b_1}{}\\
	&= (\equivrep{a_1, b_1}{} \times \equivrep{a_2,b_2}{}) + (\equivrep{a_1, b_1}{} \times  \equivrep{a_3, b_3}{})\\
	&= (i \times j) + (i \times k)
\end{align*}
\end{proof}

اثبات پنج شرط باقی‌مانده را به‌عنوان تمرین واگذار می‌کنیم. با انجام آن
نشان داده‌ایم که $\Int$ حلقه‌ای جابجایی تشکیل می‌دهد؛ یعنی جمع و ضرب
(چنان‌که تعریف شده‌اند) همان‌گونه رفتار می‌کنند که باید.

\begin{prob}
ثابت کنید $\Int$ حلقه‌ای جابجایی است.
\end{prob}

اما کار ما هنوز پایان نیافته است. افزون بر تعریف جمع و ضرب روی $\Int$،
رابطهٔ ترتیبِ $\leq$ را تعریف کردیم و باید بررسی کنیم که این رابطه چنان‌که
باید رفتار می‌کند. با جزئیات بیشتر، باید نشان دهیم $\Int$ حلقه‌ای
\emph{مرتب} تشکیل می‌دهد.\footnote{به یاد آورید که بنا بر
	\olref[sfr][rel][ord]{def:linearorder}، ترتیب تام رابطه‌ای بازتابی،
	تراگذری، پادتقارن و همبند است. در بافت روابط ترتیب، همبندی را گاهی
	\emph{سه‌شقی‌بودن} می‌نامند، زیرا برای هر $a$ و $b$ داریم $a \leq b \lor a
	= b \lor a \geq b$.}

\begin{defn}
یک \emph{حلقهٔ مرتب} حلقه‌ای جابجایی است که به رابطهٔ ترتیب تامِ $\leq$
نیز مجهز است، به‌گونه‌ای که:
\begin{align*}
	a \leq b &\lif a + c \leq b + c\\
	(a \leq b \land 0 \leq c) &\lif a \times c \leq b \times c
\end{align*}
\end{defn}

\begin{prob}
ثابت کنید $\Int$ حلقه‌ای مرتب است.
\end{prob}

مانند پیش، نشان‌دادن اینکه $\Int$، چنان‌که ساخته شده، حلقه‌ای مرتب است
پرزحمت اما عادی است. این کار را به شما واگذار می‌کنیم.

بدین‌ترتیب کار اعداد صحیح به پایان می‌رسد. اما اکنون باید مطالب بسیار
مشابهی را دربارهٔ اعداد گویا نشان دهیم. به‌ویژه، اکنون باید نشان دهیم که
اعداد گویا یک \emph{میدان} مرتب را تحت تعریف‌های داده‌شدهٔ $+$، $\times$
و $\leq$ تشکیل می‌دهند:
\begin{defn}\ollabel{orderedfield}
یک \emph{میدان مرتب} حلقه‌ای مرتب است که شرط زیر را نیز برآورده می‌کند:
\begin{align*}
	\emph{وارون ضربی}& & (\forall a \in S \setminus \{0\})(\exists b \in S) a\times b& = 1
\end{align*}
\end{defn}

پس از آنکه نشان دادید $\Int$ حلقه‌ای مرتب تشکیل می‌دهد، نشان‌دادن اینکه
$\Rat$ میدانی مرتب تشکیل می‌دهد آسان اما پرزحمت است.

\begin{prob}
ثابت کنید $\Rat$ میدانی مرتب است.
\end{prob}

پس از پرداختن به اعداد صحیح و گویا، تنها اعداد حقیقی باقی می‌مانند.
به‌ویژه، باید نشان دهیم $\Real$ یک میدان مرتبِ \emph{کامل}، یعنی میدان
مرتبی با خاصیت تمامیت، تشکیل می‌دهد. پیش‌تر \olref[cuts]{realcompleteness}
ثابت کرد که $\Real$ خاصیت تمامیت دارد. بااین‌حال، هنوز باید بررسیِ
ملال‌آورِ این مطلب را انجام دهیم که $\Real$ میدانی مرتب است.

پیش از آنکه به‌سراغ \emph{آن} تمرین پرزحمت برویم، باید چند مطلب
«بی‌واسطه‌تر» را بررسی کنیم. برای نمونه، باید تضمین کنیم که عبارت
$\alpha + \beta$، چنان‌که تعریف شده، واقعاً یک \emph{برش} است، برای هر
دو برش $\alpha$ و $\beta$. این هم برهان آن:

\begin{proof}	
چون $\alpha$ و $\beta$ هر دو برش‌اند، $\alpha + \beta = \Setabs{p
+ q}{p \in \alpha \land q \in \beta}$ زیرمجموعه‌ای ناتهی و حقیقی از
$\Rat$ است. اکنون فرض کنید $x < p + q$ برای بعضی $p \in \alpha$ و
$q \in \beta$. آنگاه $x - p < q$، پس $x - p \in \beta$، و $x = p + (x - p) \in
\alpha + \beta$. بنابراین $\alpha + \beta$ پاره‌ای آغازین از $\Rat$ است.
در پایان، برای هر $p + q \in \alpha + \beta$، چون $\alpha$ و $\beta$
هر دو برش‌اند، $p_1 \in \alpha$ و $q_1 \in \beta$ای وجود دارند چنان‌که
$p < p_1$ و $q < q_1$؛ پس $p + q < p_1 + q_1 \in \alpha +
\beta$؛ بنابراین $\alpha + \beta$ بیشینه ندارد.
\end{proof}

با تلاش‌هایی مشابه می‌توانید بررسی کنید که $\alpha - \beta$ و
$\alpha \times \beta$ و $\alpha \div \beta$ برش‌اند (در مورد آخر، با
کنارگذاشتن حالتی که $\beta$ برشِ صفر است). این بار نیز کار را به شما
واگذار می‌کنیم.

\begin{prob}
ثابت کنید $\Real$ میدانی مرتب است.
\end{prob}

اما نکتهٔ کوچکِ بازمانده‌ای هست که باید سامان دهیم. در
\olref[cuts]{sec} پیشنهاد کردیم که می‌توانیم $\sqrt{2} =
\Setabs{p \in \Rat}{p < 0 \text{ یا }p^2 < 2}$ بگیریم. ولی باید نشان
دهیم این مجموعه یک \emph{برش} است. این هم برهان آن:

\begin{proof}
آشکارا این مجموعه پارهٔ آغازینِ ناتهی و حقیقی از اعداد گویاست؛ پس کافی است
نشان دهیم بیشینه ندارد. به‌ویژه، کافی است نشان دهیم هرگاه $p$ عددی گویای
مثبت باشد که $p^2 < 2$ و $q = \frac{2p+2}{p+2}$، هم $p < q$ و هم
$q^2 < 2$. برای دیدن $p < q$، تنها توجه کنید:
\begin{align*}
	p^2 &< 2\\
	p^2 + 2p &< 2 + 2p\\
	p(p + 2) &< 2 + 2p\\
	p &< \tfrac{2+2p}{p+2} = q
\end{align*}
برای دیدن $q^2 < 2$، تنها توجه کنید:
\begin{align*}
	p^2 &< 2\\
	2p^2 + 4p + 2 &< p^2 + 4p+ 4\\
	4p^2 + 8p + 4 &< 2(p^2 + 4p + 4)\\
	(2p+2)^2 & <2(p+2)^2\\
	\tfrac{(2p+2)^2}{(p+2)^2} &< 2\\
	q^2 &< 2
\end{align*}
\end{proof}
\end{document}
